% 第 14 章 拉格朗日力学：换一种语言讲运动
\chapter{拉格朗日力学：换一种语言讲运动}

\step{反常识开场}

\begin{mainline}
科技馆里常能见到一件展品：两根等长的轻杆首尾铰接，每根杆的末端挂一个金属球，这叫双摆。把它拉到一个角度松手，接下来发生的事情完全不像钟摆那样斯文——它旋转、翻跟头、忽然加速忽然减速，轨迹像醉汉走路，永不重复。你可以回家用两把钥匙和两段线绳自己做一个：第一根线吊在门把手上，第二根线系在第一把钥匙上。松手之后盯十秒钟，你找不到任何规律；再试一次，明明觉得放手姿势一模一样，轨迹却完全不同。

再看看你的手腕。如果你戴的是自动机械表，表壳里有一枚半圆形的重金属摆陀，随手腕的动作来回转动，把发条一点点上紧。工程师设计它时面对的问题和双摆一模一样：摆陀绕着一根轴转，轴上还有轴承的约束，手腕的晃动毫无规律——他们怎么预言并优化这个小机器的运动？

这两个装置似乎是对物理学的嘲笑：第 5 章教我们用位置和速度描述运动，第 6 章教我们用 \(F=ma\) 预言运动，可是面对双摆和摆陀，谁写得出它们的方程？

令人意外的事实是：物理学家不但写得出，而且写得比单摆难不了多少。更奇怪的是，他们写方程时从头到尾都没有计算两段杆之间、杆与支架之间、摆陀与轴承之间的相互作用力——那些力是牛顿方法里最难缠的部分，而在新方法里它们干脆不出现。

不算力，凭什么能算对运动？这就是本章要回答的问题。答案是：不是牛顿错了，而是“逐点分析受力”这种语言不适合这类问题。物理学家为同一套自然规律发明了第二种语言——拉格朗日力学。学会它，你就拿到了机器人工程师、航天器设计师每天都在用的那套工具。
\end{mainline}

\step{先猜一猜}

\begin{mainline}
在往下读之前，先把你的直觉摆到桌面上，回答四个问题：

第一，描述双摆每一瞬间的状态，至少需要几个数？直觉可能说：两个球，每个球要横、纵两个坐标，共四个数。也可能你说不出理由。

第二，一粒珠子穿在光滑的弯铁丝上，从一端松手滑下。要预言它每一时刻的位置和快慢，是否必须先算出铁丝在每一个位置对珠子的支持力？多数人的直觉是：当然要，不知道力怎么算运动？

第三，还是那粒珠子。不计摩擦，它从高处滑下再冲上对面的坡。它有没有可能冲到比出发点更高的地方？请凭直觉回答，并想清楚你的理由是什么。

第四，同一个单摆，挂在静止的电梯里和挂在匀加速上升的电梯里，摆动快慢一样吗？别急着套公式，先猜方向：变快、变慢，还是不变？

记住你的四个答案。本章结束时你会发现：第一题的答案是两个数，不是四个；第二题的答案是完全不需要；第三题你的理由里藏着整章的钥匙；第四题则会免费送给你一个理解失重的角度。
\end{mainline}

\step{动手实验}

\begin{experiment}[衣架轨道与珠子]
找一根可以弯折的粗铁丝（把金属晾衣架拉直就行），再找一个能穿在铁丝上的东西：一颗带孔的算盘珠、一个钥匙圈、甚至一个垫片。把铁丝弯成一道“先下坡、后上坡”的滑道，两端固定在不同高度，出口端比入口端略低一点，把珠子穿上去。

第一步，把珠子放在高端松手。观察它滑到对面坡上的最高点：是不是几乎回到出发时的高度，但总差一点点？多试几次，差的那一点点很稳定。

第二步，把滑道中间某一段掰得更陡。再松手，你会发现珠子在陡坡上加速明显更猛，在对面上坡时却照样停在老地方附近——它从来冲不过出发点的高度。

第三步，用手机慢动作录像。逐帧看，珠子在低处飞快、在高处磨蹭。快和慢不是乱来的：高度降多少，速度就长多少，像有本账在严格兑换。

如果你还有耐心，把两颗珠子同时从不同高度松手，比赛谁先到对面。你会发现一个起初违反直觉的结果：从更低处出发的珠子，未必更慢——它在低处已经攒够了速度。这条赛道，赢家取决于整本账，而不是某一瞬间的快慢。

最后试着定性回答：为什么珠子永远冲不过出发点？别翻书，用你自己的话把理由说出来，再和同伴争论一轮。这个“不许超过”的限制，就是本章后面一切方程的源头之一。
\end{experiment}

\begin{mainline}
这个实验值得多看一眼的地方在于：想用力学原理解释它，你会立刻感到别扭。铁丝对珠子的支持力时刻在改变方向——轨道每弯一次，方向就变一次；它的大小也不知道，得先知道运动才能算力，可要算运动又得先知道力。一个死循环。

但换一种说法，事情简单得不像话：高度换来的速度，沿轨道一分不多一分不少；能量这本账，摩擦每漏掉一点，最高点就矮一点。第 9 章的能量守恒在这里运转良好，而受力分析在这里寸步难行。这个对比不是巧合，它在提示：也许存在一种语言，从能量出发、从“你选的变量”出发，而不是从力出发。本章的正式内容，就是把这种语言造出来。

顺便回应全书反复出现的三个追问。系统此刻是什么状态？在这个实验里，珠子在轨道上的位置加它的快慢，两个数就够了。什么规律决定状态如何变化？暂时是能量这本账，本章会把它升级成更强大的规则。我们怎样通过测量知道它？慢动作录像加尺子，就是你手边的全部仪器。
\end{mainline}

\step{旧模型遇到麻烦}

\begin{mainline}
先郑重声明：牛顿第二定律没有任何错。直到今天，发射火箭、设计桥梁，底层的物理仍然是它。麻烦出在表达上，具体说有三层。

第一层麻烦：约束力是“事后才知道”的力。单摆的绳张力、铁丝对珠子的支持力、双摆铰链处的相互作用，都不是事先写在题目里的已知量，而是要靠方程解出来的未知数。它们的方向还随运动不断变化。在第 7 章处理斜面时这已经让人头疼，那时约束力至少是常数；到了双摆，约束力本身就像个醉汉。

第二层麻烦：默认的坐标不合适。用横坐标、纵坐标描述单摆，需要两个数，外加一句“它离悬点的距离永远是绳长”——一句写在方程之外的附加条件。可谁都看得出来：一个角度就够了。我们在用两个变量加一条补充说明，描述一个本质上只有一个自由度的系统，多出来的东西全是自己给自己找的麻烦。

第三层麻烦：系统一复杂，方程数量爆炸。拿双摆认真数一数：两个球各要横、纵两个坐标，共四个；两条“杆长不变”的约束，是两条额外的方程；铰链和支架处的约束力共四个未知分量。于是你面对八个未知数、八条互相纠缠的方程，而其中一半的未知数你根本不关心——没人想知道铰链此刻的受力，大家只想知道球下一瞬在哪里。换成六个关节的机械臂，未知数的队伍排得更长，而工厂里的控制器每秒要把它们解上千遍。

打个比方。你在一个陌生城市打车，一种说法是用经纬度报目的地，另一种说法是“前面路口左转，直行过两个红绿灯，右手边就是”。它们像什么？像两种语言描述同一座城市，指的是同一个地方。它们不像什么？不像两张互相矛盾的地图——如果两种说法把你带到了不同的地方，那才是出了真问题。牛顿语言和本章要发明的语言也是这个关系：它们是同一套自然规律的两种说法，给出的预言必须处处一致。我们要换语言，不是因为旧语言错，而是因为旧语言在这类问题上太繁琐。
\end{mainline}

\step{发明新概念}

\begin{mainline}
现在像历史上的物理学家一样，一步步把新语言造出来。拉格朗日在 1788 年出版《分析力学》时颇为自豪地说：这本书里没有一幅图，也没有一次受力分析。他的办法可以拆成三个发明。

\textbf{发明一：广义坐标。}先数一数系统真正“自由”的个数，叫自由度。单摆：一个角度 \(\theta\)，自由度一。铁丝上的珠子：沿轨道的弧长 \(s\)，自由度一。双摆：两个角度，自由度二。六关节机械臂：六个关节角，自由度六。然后放手去选最顺手的变量，并给它们一个统一的名字 \(q\)。你不再被“横坐标纵坐标”绑架：旋转木马上的位置用角度，弹簧振子用伸长量，活塞用推进深度，什么方便用什么。

广义坐标甚至可以不唯一。描述操场上的同学，你可以用“离旗杆多远、朝哪个方向”，也可以用“沿跑道跑了多少米”——只要两个数能唯一确定位置，就是好坐标。选哪一套，全看哪套让问题简单。这是一种前所未有的自由：变量是你为问题量身定做的，而不是问题硬塞给你的。

\textbf{发明二：用广义坐标写两本账。}一本是动能 \(T\)，记录“此刻动得有多猛”；一本是势能 \(V\)，记录“此刻位置值多少”。第 9 章说过它们像能量的两个口袋。关键在于：这两本账用你自己选的 \(q\) 来写。单摆的动能写成 \(T=\frac{1}{2}ml^2\dot{\theta}^2\)，势能以最低点为零点写成 \(V=-mgl\cos\theta\)，绳张力从头到尾没有登场的机会。

为什么约束力可以这样被晾在一边？秘密在坐标的选择里。珠子的坐标 \(s\) 只能沿轨道取值，也就是说，我们发明的语言里根本不存在“离开轨道”这种说法。约束力唯一的工作就是把珠子按在轨道上；而当你选用的坐标只允许沿轨道的运动时，这份工作已经由坐标本身默默完成了。约束力不是被近似掉，而是被翻译成了坐标系的形状。

\textbf{发明三：拉格朗日量。}定义一个新对象
\[
L=T-V,
\]
动能减势能。注意，不是加！动能加势能是能量，那是个守恒量；动能减势能是另一回事，它一般不守恒，也不是你能拿仪器直接测出来的东西。为什么要发明这么一个古怪的组合？第 13 章已经埋过伏笔：把 \(L\) 沿整条运动路径随时间累计起来，得到的总账叫作用量。真实的运动有一个惊人的性质——在所有从同一起点到同一终点的假想路径里，真实路径让这本总账停在“平台”上：轻轻改动路径，总账在一阶近似下不变。

拉格朗日量像什么？像一张每瞬间刷新的“收支差流水单”：动能是进账，势能是欠账，它记的是两者之差。它不像什么？它不像能量那样是个可以被守恒定律看管、被仪表测出的物理量。它是记账用的中间工具，物理内容不在它单独一行的数值里，而在它沿整条路径的总账里。把工具当成实物，是本章最常见的误解，后面还会专门清算它。

说一句历史，免得你以为这套语言是凭空掉下来的。十八世纪的数学家们早就被一类问题吸引：光走最短时间的路（第 13 章），绳子垂成最省势能的形状，肥皂膜绷成面积最小的曲面——“自然懂得挑最优”的想法在空气里飘了几十年。莫佩尔蒂把它写成一条含糊的“最小作用量原理”，欧拉把它在单质点情形里算严格了，拉格朗日则完成了关键一跃：他证明这条原理不是对牛顿定律的补充，而是能把全部牛顿力学重新推导出来的另一套地基。又过了几十年，哈密顿把作用量改写成今天教科书里的样子。所以你手里这套语言，是一百多年接力的成品，不是某位天才一夜的灵感。

也值得说清楚这套语言里“解释”站在哪一层。实验事实是：双摆那样摆、珠子冲不过出发点。数学模型是：作用量取平台值加欧拉—拉格朗日方程。至于“自然是不是真的有意识地比较过所有路径”——那不是实验能回答的问题，只是一种说法；有的物理学家喜欢它，有的认为它只是账本的拟人化。本书站在谨慎的一边：我们断言账本有效，不断言自然有个会计。
\end{mainline}

\begin{figure}[htbp]
\centering
\begin{tikzpicture}[>=Stealth,scale=1.1]
  % 支点
  \fill[pattern=north east lines] (-1.2,0) rectangle (1.2,0.25);
  \draw[thick] (-1.2,0) -- (1.2,0);
  \fill (0,0) circle (1.5pt);
  % 竖直虚线
  \draw[dashed] (0,0) -- (0,-3.2);
  % 绳与摆球
  \draw[thick] (0,0) -- (50:3);
  \fill (50:3) circle (4pt);
  % 角度
  \draw[->] (0,-1.3) arc (-90:50:1.3);
  \node at (-0.62,-1.55) {$\theta$};
  \node[right] at (50:3.35) {$m$};
  \node[left] at ($(0,0)!0.55!(50:3)$) {$l$};
\end{tikzpicture}
\caption{单摆：横、纵两个坐标加一条“绳长不变”的约束，不如一个广义坐标 \(\theta\) 干净。}
\end{figure}

\step{一句话模型}

\begin{oneline}
选好你顺手的坐标，写下动能减势能这本账；真实发生的运动，是使这本账沿路径的累计取“平台值”的那条路——约束力不做功时，它会在这套语言里自动消失。
\end{oneline}

\begin{mainline}
这一句话里有三个词值得停一停。“选好你顺手的坐标”是自由度的自由；“取平台值”是第 13 章最小作用量思想的重申——记住是“平台”而不是“最小”，就像站在马鞍的中心，沿某些方向是最低、沿另一些方向是最高；“自动消失”则是本章最实际的收获：支持力、张力这类约束力，只要它们不对系统做功，就根本不必出场。一句话装不下全部细节，但它抓住了灵魂：换坐标、记两账、取平台。下一节把这句话翻译成可以动手算的方程。
\end{mainline}

\step{数学翻译器}

\begin{mainline}
数学翻译器只有一条规则，叫欧拉—拉格朗日方程。对每一个广义坐标 \(q\)，做三步操作：先拿 \(L\) 对速度 \(\dot{q}\) 求“偏变化率”，再把它对时间求变化率，然后减去 \(L\) 对位置 \(q\) 的偏变化率，令结果等于零：
\[
\frac{\mathrm{d}}{\mathrm{d}t}\frac{\partial L}{\partial \dot{q}}-\frac{\partial L}{\partial q}=0 .
\]
看起来吓人，用起来是流水线。先用全世界最简单的问题试刀：一个只受重力的落体，坐标取高度 \(y\)。动能 \(\frac{1}{2}m\dot{y}^2\)，势能 \(mgy\)，于是 \(L=\frac{1}{2}m\dot{y}^2-mgy\)。\(L\) 对 \(\dot{y}\) 的偏变化率是 \(m\dot{y}\)，对时间再求一次是 \(m\ddot{y}\)；\(L\) 对 \(y\) 的偏变化率是 \(mg\)。方程给出 \(m\ddot{y}+mg=0\)，即 \(\ddot{y}=-g\)。质量照例被约掉——第 6 章说过，这正是引力质量与惯性质量相等的表现，新语言原样继承了这条深刻的平等。

再看一个老朋友。弹簧振子：动能 \(T=\frac{1}{2}m\dot{x}^2\)，势能 \(V=\frac{1}{2}kx^2\)，于是 \(L=\frac{1}{2}m\dot{x}^2-\frac{1}{2}kx^2\)。第一步，\(L\) 对 \(\dot{x}\) 的偏变化率是 \(m\dot{x}\)，对时间再求一次是 \(m\ddot{x}\)；第二步，\(L\) 对 \(x\) 的偏变化率是 \(-kx\)。相减为零：\(m\ddot{x}+kx=0\)，正是 \(F=-kx\) 的改写。新语言在简单问题上先交了一份和旧答案完全一致的答卷——这是它必须通过的入职考试。

来一个约束力真正登场的例子：阿特伍德机，即滑轮两侧各挂一个砝码，质量 \(m_1\) 和 \(m_2\)。牛顿做法要设绳张力为未知数、对两个砝码各列方程再消元。新语言里，绳长不变意味着“一边下降多少、另一边就上升多少”，所以整个系统只有一个自由度：取 \(m_1\) 下降的距离 \(x\)。动能是 \(\frac{1}{2}(m_1+m_2)\dot{x}^2\)，势能随 \(x\) 的变化是 \((m_2-m_1)gx\)，方程一步给出
\[
\ddot{x}=\frac{m_1-m_2}{m_1+m_2}\,g .
\]
立刻检查：两边等重，\(\ddot{x}=0\)，系统平衡或匀速滑动，合理；\(m_2=0\)，\(\ddot{x}=g\)，\(m_1\) 自由落体，合理；加速度永远小于 \(g\)，因为另一边的砝码在拖后腿，合理。张力？一次都没出现，但如果你想知道它，把解出的 \(\ddot{x}\) 代回任一边的牛顿方程即可反推。

再看单摆：\(L=\frac{1}{2}ml^2\dot{\theta}^2+mgl\cos\theta\)。第一步，\(L\) 对 \(\dot{\theta}\) 求偏变化率得 \(ml^2\dot{\theta}\)；对时间再求一次得 \(ml^2\ddot{\theta}\)。第二步，\(L\) 对 \(\theta\) 求偏变化率得 \(-mgl\sin\theta\)。令两者相减为零：
\[
ml^2\ddot{\theta}+mgl\sin\theta=0
\quad\Longrightarrow\quad
\ddot{\theta}=-\frac{g}{l}\sin\theta .
\]
这正是第 6 章用受力分析推出的单摆方程。两套语言，同一个答案——这正是上一节“两种语言同一座城市”的兑现。

按老规矩，立刻用特殊情况检查这个公式。令 \(g=0\)：\(\ddot{\theta}=0\)，空间站里悬挂的小球不再摆动，角速度是零就永远停着，合理。把 \(l\) 取得极大：摆动极慢，合理的——井口垂下的长绳半天才晃一个来回。令 \(\theta\) 极小：\(\sin\theta\approx\theta\)，方程变成 \(\ddot{\theta}=-(g/l)\theta\)，这是第 16 章要细讲的简谐振动，周期 \(T=2\pi\sqrt{l/g}\)。

顺手做一次带真实数据的估算。公园秋千的绳长约 \(2\) 米，代入 \(g=9.8\ \mathrm{m/s^2}\)：
\[
T=2\pi\sqrt{\frac{2}{9.8}}\approx 2.8\ \mathrm{s}.
\]
下次荡秋千数一下：从最高点回到同一侧最高点，大约两三个心跳，正是三秒上下。注意这个周期里既没有你的体重，也没有绳的张力——质量 \(m\) 在推导中被约掉了，张力压根没出现过。
\end{mainline}

\begin{mathbridge}
为什么“令账本的偏变化率为零”就是运动定律？补上第 13 章留下的推导。把作用量写成
\[
S=\int_{t_1}^{t_2} L\bigl(q,\dot{q}\bigr)\,\mathrm{d}t .
\]
假设真实路径是 \(q(t)\)，给它加一个微小扰动：\(q(t)\to q(t)+\varepsilon\eta(t)\)，其中 \(\eta\) 在两端为零（起点终点不变）。作用量的一阶改变量是
\[
\delta S=\int_{t_1}^{t_2}\left(\frac{\partial L}{\partial q}\eta+\frac{\partial L}{\partial \dot q}\dot{\eta}\right)\mathrm{d}t .
\]
对第二项做分部积分，端点项因 \(\eta\) 在两端为零而消失：
\[
\delta S=\int_{t_1}^{t_2}\left(\frac{\partial L}{\partial q}-\frac{\mathrm{d}}{\mathrm{d}t}\frac{\partial L}{\partial \dot q}\right)\eta\,\mathrm{d}t .
\]
“平台”的意思是：对任意扰动 \(\eta\)，都有 \(\delta S=0\)。要让任意形状的 \(\eta\) 乘上一个函数再积分都为零，唯一的可能是括号里处处为零——这正是欧拉—拉格朗日方程。整个推导只做了一件事：把“账本取平台值”这句话翻译成一个坐标一个方程。

顺带认识一个以后会反复出现的角色。\(L\) 对 \(\dot{q}\) 的偏变化率 \(p=\partial L/\partial\dot{q}\) 叫广义动量：对弹簧振子它是 \(m\dot{x}\)，就是普通动量；对单摆它是 \(ml^2\dot{\theta}\)，正是角动量。欧拉—拉格朗日方程于是可以读成：广义动量的变化率等于 \(L\) 对位置的偏变化率——这正是“动量变化率等于力”的广义版本。第 15 章的哈密顿力学，就是把这对搭档 \((q,p)\) 扶正当主角。
\end{mathbridge}

\begin{figure}[htbp]
\centering
\begin{tikzpicture}[>=Stealth,scale=1.0]
  % 坐标轴
  \draw[->] (0,0) -- (6.4,0) node[below] {$t$};
  \draw[->] (0,0) -- (0,3.2) node[left] {$q$};
  % 端点
  \fill (0.8,0.7) circle (2pt) node[below left] {起点};
  \fill (5.6,2.4) circle (2pt) node[above right] {终点};
  % 真实路径
  \draw[thick] (0.8,0.7) .. controls (2.2,1.0) and (4.2,2.2) .. (5.6,2.4);
  % 假想路径
  \draw[thick,dashed] (0.8,0.7) .. controls (2.0,2.2) and (4.0,0.9) .. (5.6,2.4);
  \node at (3.1,1.75) {真实路径};
  \node at (3.0,0.62) {假想路径};
\end{tikzpicture}
\caption{固定起点与终点，在所有假想路径中，真实路径使作用量取“平台值”：轻轻一推路径，账单一阶不变。}
\end{figure}

\begin{challenge}
感受这套语言真正的威力，看一个受力分析极易做错的系统：光滑水平面上放一个质量 \(M\) 的斜面，斜面上放一个质量 \(m\) 的小滑块，斜面倾角 \(\alpha\)，一切接触面光滑。松手后滑块下滑，斜面同时被顶得向后退——直觉常在这里翻车：很多人下意识把斜面当成不动的，直接写 \(a=g\sin\alpha\)。

用拉格朗日语言：取斜面位置 \(X\) 和滑块沿斜面下滑的距离 \(s\) 两个广义坐标。滑块相对地面的速度是“斜面速度”与“沿斜面速度”的矢量合成，动能为
\[
T=\tfrac12 M\dot X^2+\tfrac12 m\left(\dot X^2+2\dot X\dot s\cos\alpha+\dot s^2\right),
\qquad
V=-mgs\sin\alpha .
\]
对 \(X\) 和 \(s\) 各写一条欧拉—拉格朗日方程，联立解出
\[
\ddot s=\frac{(M+m)\,g\sin\alpha}{M+m\sin^2\alpha},
\qquad
\ddot X=-\frac{m\,g\sin\alpha\cos\alpha}{M+m\sin^2\alpha}.
\]
检查极限：令 \(M\to\infty\)，斜面重如大地，\(\ddot s\to g\sin\alpha\)、\(\ddot X\to 0\)，回到固定斜面的老答案；令 \(\alpha=0\)，一切为零，平面上什么也不发生；令 \(\alpha\to 90^\circ\)，滑块几乎自由落体，\(\ddot s\to g\)，斜面纹丝不动。全部通过。注意全程没有出现滑块与斜面之间的支持力——它被坐标的选择消掉了，而它被消掉得\textbf{精确}，不是近似。另外，\(L\) 中不含 \(X\) 本身（只含 \(\dot X\)），于是 \(X\) 方向的广义动量守恒——这正是水平方向总动量守恒，第 10 章的老朋友以新面孔出现了。
\end{challenge}

\step{模型的边界}

\begin{mainline}
拉格朗日力学不是万能钥匙，它的地盘有清楚的边界。

它擅长的场合：约束是“理想”的——约束力始终垂直于允许的运动方向，因而对系统不做功。光滑轨道、刚性杆、绷紧不可伸长的绳都属于这类。这时约束力精确地从方程中消失。它要求主动力能用势能写出，也就是第 9 章说的保守力：重力、弹力、引力都行。

它需要修补的场合：摩擦和空气阻力。这些力做负功、没有势能可写，只能作为“广义力”额外塞回方程右边，或者引入耗散函数。方程仍然成立，但优雅程度打折。这是实话：碰到满身摩擦的问题，新语言的省力程度会缩水。还有一类更微妙的麻烦：有些约束没法写成“坐标之间的等式”，比如在地面上纯滚动的球，它的位置和转角被一条涉及速度的关系捆在一起。这类问题需要额外的技术（拉格朗日乘子），本章不展开，但你至少应该知道：并不是所有“被捆住”的系统都能用数自由度的办法一口吞下。

还有一个容易踩空的台阶：守恒。什么时候 \(T+V\) 这本总账守恒？答案是：当 \(L\) 的表达式里不显含时间、约束也不随时间变化时。荡秋千的人在最高点下蹲、最低点站起，等于一边荡一边改动摆长，账本的规则本身在变，机械能不再守恒——多出来的能量是肌肉化学能付的账。这不是定律失败，而是适用条件被悄悄改掉了。第 15 章会把“什么量守恒”这个问题变成一套系统的问法。

它彻底让位的场合：尺度小到量子领域，粒子根本没有唯一确定的轨迹，“从所有路径里挑一条”必须让位给“所有路径一起贡献”——第 43 章会讲量子态，第 50 章会讲费曼如何把这个想法变成量子力学的另一种表述。有趣的是，费曼的出发点恰恰就是本章的账本：每条路径都按作用量领一个相位。经典力学的“挑一条路”在那里变成了量子规则的近似。

再清算三个典型误用。第一，把 \(L=T-V\) 当成能量。能量是 \(T+V\)，两者符号差一处，含义完全不同；把记账工具当成实物，后面的账全会错。第二，以为拉格朗日力学和牛顿力学是两派互相竞争的自然理论，仿佛实验要在二者之间投票。不是。实验事实只有一个——双摆怎么摆就怎么摆；两种力学是描述同一事实的两种数学语言，预言必须一致，不一致的那次一定是有人算错了。第三，以为作用量永远取“最小”。它取的是驻值即平台值：第 13 章讲光路时说过，反射、折射给出的常常是极大或马鞍点，力学路径同理。

最后是一个会让直觉出错的反例。不少人第一次学到这里会想：不计算约束力，是不是丢了信息、只是个近似解法？恰恰相反。只要约束理想，约束力被消元是严格的，解出的运动与“老老实实解出约束力再积分”的结果一模一样，连小数点后面的每一位都相同。如果你真想知道绳张力多大（比如校核绳子会不会断），可以先把运动解出来，再回头反推约束力——它并没有丢，只是不必在解题的路上陪跑。
\end{mainline}

\step{回头解释开场现象}

\begin{mainline}
回到科技馆那个疯狂的双摆。现在你知道该怎么对付它了。

第一步，选广义坐标：第一根杆与竖直方向的夹角 \(\theta_1\)，第二根杆与竖直方向的夹角 \(\theta_2\)。两个数，完整描述整个系统的任意瞬间——这就是开场第一题的答案。

第二步，写两本账。动能是两球动能之和：第一个球好办，第二个球的速度带着第一根杆的牵连，展开后含一项把 \(\theta_1\) 和 \(\theta_2\) 耦合在一起的交叉项；势能是两球高度之和。写下来，\(L\) 只有一行多长。

第三步，走两遍欧拉—拉格朗日流水线，得到两个耦合的二阶方程。没有解析解——这不是方法失败，而是这个系统本身如此，混沌意味着任何方法都写不出封闭公式。但交给计算机逐步积分，屏幕上重现的正是你在科技馆看到的醉汉舞蹈：旋转、翻跟头、对初始条件的极端敏感——开头差一根头发丝的角度，几秒后轨迹就南辕北辙，这就是你两次放手得到不同舞蹈的原因。两段杆之间、杆与支架之间的相互作用力，从头到尾没有出现。它们不是被忽略了，而是被“选对坐标”这个动作精确地消了元。

手腕上的摆陀同理：一个转角、一本动能账、一本发条的势能账，方程就有了，轴承力不必出场。工程师剩下的工作是更有意思的部分——调整摆陀的质量分布，让账本在日常手腕动作下给发条上最多的链。

这套流程不是书斋里的游戏。工业机器人的控制器用它工作：一条六关节机械臂，用牛顿语言要处理十八个坐标加十几条约束和一大堆未知约束力，用拉格朗日语言只有六个方程，而控制器每秒要解上千遍——省下的计算量就是机械臂能实时抓取零件的原因之一。汽车悬架的多体仿真、火箭姿态控制、动画电影里角色的物理引擎，底层都是同一套语言。开场的双摆和工厂里的机械臂，共享同一本账。

开场第四题也顺手结了账：电梯匀加速上升时，电梯里的一切像处在等效重力 \(g+a\) 中，单摆周期变成 \(2\pi\sqrt{l/(g+a)}\)，变短了；钢缆断掉的瞬间，等效重力归零，摆停在半空——和第 8 章的失重是同一回事。你没有学新定律，只是旧账本换了参照系重新记了一遍。
\end{mainline}

\begin{figure}[htbp]
\centering
\begin{tikzpicture}[>=Stealth,scale=1.0]
  % 支架
  \fill[pattern=north east lines] (-1.0,0) rectangle (1.0,0.22);
  \draw[thick] (-1.0,0) -- (1.0,0);
  \fill (0,0) circle (1.5pt);
  % 第一段杆
  \draw[thick] (0,0) -- (-35:2.2);
  \fill (-35:2.2) circle (3.5pt);
  % 第二段杆（从第一球端点出发）
  \coordinate (p1) at (-35:2.2);
  \draw[thick] (p1) -- ++(-80:1.8);
  \fill ($(p1)+(-80:1.8)$) circle (3.5pt);
  % 角度标注
  \draw[dashed] (0,0) -- (0,-1.6);
  \draw[->] (0,-1.0) arc (-90:-35:1.0);
  \node at (-0.52,-1.05) {$\theta_1$};
  \draw[dashed] (p1) -- ++(0,-1.2);
  \draw[->] ($(p1)+(0,-0.75)$) arc (-90:-80:0.75);
  \node[right] at ($(p1)+(0.06,-0.8)$) {$\theta_2$};
  \node[left] at ($(0,0)!0.5!(p1)$) {$l_1$};
  \node[right] at ($(p1)!0.5!($(p1)+(-80:1.8)$)$) {$l_2$};
\end{tikzpicture}
\caption{双摆只需两个广义坐标 \(\theta_1\)、\(\theta_2\)；铰链处的相互作用力无需出场。}
\end{figure}

\step{脑洞实验与挑战题}

\begin{thought}[把行星装进账本，再把账本推向极端]
这套语言能撑到多大的尺度？把地球绕太阳的运动写进来：广义坐标取距离 \(r\) 和方位角，动能是径向与横向两部分之和，势能是 \(-GMm/r\)，两条欧拉—拉格朗日方程推出来的正是第 8 章见过的行星运动，还能白捡一个守恒量——方位角不出现在 \(L\) 里，它对应的广义动量（就是角动量）自动守恒。再贪心一点：把太阳系的行星全写进同一个 \(L\)，每个行星两三个坐标，账本变厚，规则不变。

现在推向极端。如果系统的“坐标”不是几个数，而是空间里每一点上的一个值——比如鼓面上每一点的高度、电磁场在每一点的强度——那相当于有无穷多个广义坐标。拉格朗日的办法居然照样走得通：把求和换成积分，同样的“账本取平台值”原则，推出的是场的方程。第 37 章光学的统一图景和第 59 章的规范理论，都是沿着这条路走下去的。一个为机械装置发明的记账法，最后成了整个近代物理的书写格式。还有一条更近的岔路：第 15 章会把“位置加速度”这对搭档换成“位置加动量”，把这本账改写成在状态空间里航行——那是哈密顿的语言，也是通往量子力学的桥。
\end{thought}

\begin{mainline}
最后留四道题。它们都不要求你算到底，只要求你把本章的思路用在新处境里。
\end{mainline}

\begin{puzzles}
\item 铁丝上的珠子滑过一道“先下后上”的轨道。现在把轨道改成一个倒扣的圆弧顶：珠子从顶端附近出发。它有没有可能在半途离开轨道飞出去？如果会，“约束力不做功”这句话还成立吗？拉格朗日方程会在哪里发出警报？（思路提示：铁丝只能从外侧推珠子，不能从内侧拉它——约束是单向的。当你解出的运动要求一个“把珠子往回吸”的支持力时，就是约束失效、珠子飞离的时刻。方程本身不会报警，你得盯着被消元的那个力。）
\item 荡秋千时，在最高点下蹲、最低点站起，秋千会越荡越高。用本章的语言说：是谁在给你的机械能做功？\(L\) 的表达式里发生了什么变化，使得“能量守恒”这个老裁判被合法地请下了场？（思路提示：站起和下蹲改变了等效摆长，相当于账本的规则随时间变化；回想“\(T+V\) 守恒”需要什么条件，再想想肌肉里的化学能去了哪里。第 16 章会把这种周期性改变参数的做法叫参数激励。）
\item 电梯里挂一个单摆，电梯以加速度 \(a\) 匀加速上升。用本章的语言预言：摆动周期变长、变短还是不变？如果钢缆断了电梯自由下落呢？（思路提示：不必写方程，想想电梯里的一切像处在等效的重力 \(g+a\) 中；再对照第 8 章失重的讨论，自由下落时等效重力是多少，周期意味着什么。）
\item 计算机模拟双摆一小时，发现总能量缓慢增长，越涨越快。这是方程错了，还是程序错了，还是双摆真的在凭空制造能量？你打算用什么办法区分这三种可能？（思路提示：能量在这套语言里是守恒量，正好可以当“监考老师”；数值积分每步都有微小误差，想想什么情况下误差会滚雪球，以及把步长减半后结果怎么变能帮你破案。）
\end{puzzles}
